% Deutsche Parallelfassung zu foundations/state_model.tex

\chapter{Zustandsübergang}
\label{chap:state-model}

Ein Ingenieur bearbeite die Krümmung auf einem Intervall
\(J\subseteq I\), indem er \(\kappa(s)\) durch
\(\widetilde\kappa(s)=\kappa(s)+\Delta\kappa(s)\) ersetzt. Die Krümmung
ändert sich unmittelbar. Nachfolgender Richtungswinkel und Position ändern
sich durch Integration; die Überhöhung ändert sich nur, wenn ihre eigene
Steuerung bearbeitet wird oder eine Kopplungsbedingung dies erfordert. Die
Alignment-Identität, die für die Rekonstruktion gewählte Anfangsbedingung und
der anwendbare Realization Context bleiben fest, sofern die Ingenieuraufgabe
sie nicht ausdrücklich ändert.

Diese einzelne Bearbeitung begründet den Zweck eines Zustandsmodells: Es
beantwortet, welcher Zustand aus deklarierten Eingaben unter benannten
Qualifikationen folgt. Es beschreibt keinen Speicherschnappschuss und ist kein
allgemeiner Container für jede dem Alignment zugeordnete Eigenschaft.

\section{Der Zustand, der die Frage beantwortet}

\begin{definition}[Engineering State]
\label{def:state}
Ein Engineering State ist eine strukturierte Beschreibung eines Engineering
Objects unter einem bestimmten Zustand und innerhalb benannter
Qualifikationen.
\end{definition}

Für die vorliegende Frage lautet eine Zustandstrajektorie
\[
x:I\rightarrow\mathcal X,
\qquad s\longmapsto x(s).
\]
Die Abbildung hält die Größen fest, die zur Fortpflanzung der Änderung
benötigt werden. Eine Datei, eine zwischengespeicherte Polyline oder eine
Datenbankzeile kann Werte dieser Trajektorie repräsentieren; die Speicherung
bestimmt aber nicht, welchen Zustand oder welche Qualifikationen diese Werte
beschreiben.

\section{Durch Krümmung getriebener ebener Übergang}

\begin{definition}[pose2]
\label{def:notation-pose2}
Der ebene Alignment-Zustand ist
\begin{equation}
\mathrm{pose2}(s)=\bigl(\gamma(s),\theta(s)\bigr),
\label{eq:pose2-definition}
\end{equation}
wobei \(\gamma(s)\) die ebene Position und \(\theta(s)\) den
Richtungswinkel bezeichnet.
\end{definition}

Seine Entwicklung lautet
\begin{equation}
\begin{aligned}
\gamma'(s)&=
\begin{pmatrix}\cos\theta(s)\\\sin\theta(s)\end{pmatrix},\\
\theta'(s)&=\kappa(s).
\end{aligned}
\label{eq:pose2-evolution}
\end{equation}

Bei festen \(\gamma(s_0)\) und \(\theta(s_0)\) erzeugt die Anwendung von
\(\widetilde\kappa\) eine neue Zustandstrajektorie
\(\widetilde{\mathrm{pose2}}\). Innerhalb und nach dem bearbeiteten Intervall
müssen der Richtungswinkel erneut integriert und die Position rekonstruiert
werden. Eine aus der früheren Trajektorie abgeleitete bestehende CAD-Kurve
oder abgetastete Polyline ist folglich veraltet, bis ihre Repräsentation neu
erzeugt wird.

\paragraph{Ausführungsvertrag.}
Der Übergang empfängt die Krümmungsänderung, das Rekonstruktionsintervall und
die Anfangspose; er erzeugt eine ebene Zustandstrajektorie. Er setzt
Bogenlängenparametrisierung und die benannten Randbedingungen voraus. Er
bewahrt die Alignment-Identität. Sein Ergebnis erbt die Provenienz der
Änderung und der Randbedingungen. Er kann weder ableiten, ob die Änderung
zulässig ist, noch ob eine Vermessung denselben Zustand betrifft oder das
Ergebnis angenommen werden soll.

\section{Räumlicher Übergang und Steuerungen}

\begin{definition}[Grundlegender pose3-Zustand]
\label{def:pose3}
Der grundlegende räumliche Eisenbahnzustand ist
\begin{equation}
\mathrm{pose3}(s)=\bigl(\gamma(s),\theta(s),\phi(s)\bigr),
\label{eq:pose3-definition}
\end{equation}
wobei \(\phi(s)\) die Überhöhungsdrehung um die lokale Tangente ist.
\end{definition}

Dies ist die intrinsische Eingabe einer späteren räumlichen Konstruktion,
keine World-Space-Pose. Profil, metrischer Kontext und Frame-Konstruktion sind
weiterhin erforderlich, bevor sie räumlich messbar wird.

\begin{definition}[Grundlegende Steuerungen]
\label{def:control-quantities}
Das grundlegende Steuerungspaar lautet
\begin{equation}
u(s)=\bigl(\kappa(s),\omega(s)\bigr),
\qquad \omega(s)=\phi'(s).
\label{eq:pose3-controls}
\end{equation}
\end{definition}

Die Steuerungen pflanzen den Zustand fort durch
\begin{equation}
\begin{aligned}
\gamma'(s)&=
\begin{pmatrix}\cos\theta(s)\\\sin\theta(s)\end{pmatrix},\\
\theta'(s)&=\kappa(s),\\
\phi'(s)&=\omega(s).
\end{aligned}
\label{eq:pose3-evolution}
\end{equation}

\section{Freiheitsgrade während der Änderung}

Die Änderung ist nur interpretierbar, wenn der Status jeder Größe
ausdrücklich ist.

\begin{description}
\item[Fest:] Alignment-Identität, anwendbare Qualifikationen,
  Rekonstruktionsursprung und alle Randwerte, die die Aufgabe als invariant
  deklariert.
\item[Bearbeitet:] \(\kappa\) auf \(J\), mit Quelle, Urheber und Umfang als
  Provenienz.
\item[Abgeleitet:] Richtungswinkel, Tangente, ebene Position und jede aus der
  geänderten Trajektorie abgetastete Repräsentation.
\item[Gekoppelt:] Überhöhung, Profil oder Größen angrenzender Übergänge, wenn
  eine anwendbare Ingenieurbeziehung sie mit der Krümmungsänderung verbindet.
\item[Beschränkt:] Stetigkeit, Regularität, Randbedingungen, Zulässigkeit und
  projektspezifische Anforderungen aus dem anwendbaren Wissen.
\item[Frei:] Größen, die durch das deklarierte Problem weder festgelegt noch
  abgeleitet, gekoppelt oder beschränkt sind. Sie bleiben mögliche
  Freiheitsgrade; ein Solver darf sie untersuchen, kann ihnen aber keine
  Ingenieurautorität zuweisen.
\end{description}

Diese Klassifikation bildet die Brücke von der formalen Zustandsfortpflanzung
zum Entwurf von Alignment-Übergängen: Sie bestimmt die Variablen, die einem
späteren Konstruktions- oder Optimierungsverfahren zur Verfügung stehen, ohne
Software oder einen bestimmten Solver vorzuschreiben.

\section{Was verglichen werden darf}

Metrische Realisierung kann den überarbeiteten Zustand räumlich messbar machen.
Eine Vermessungspolyline kann dann Beobachtungsevidenz unter kompatiblen
Qualifikationen liefern, und ein Bewertungsoperator kann Residuen berechnen.
Der Vergleich setzt qualifizierte Zustände und Quellen in Beziehung; er macht
weder die Polyline zur physischen Realität, noch ersetzt er die konstruktive
Definition des Alignments oder verwandelt eine Bewertung in eine Entscheidung.

Das nächste Kapitel macht die Handlungen dieser Kette als typisierte
Operatoren ausdrücklich.
